\chapter{Discussion and Limitations}
\label{ch:discussion}

This chapter interprets the results of Chapter~\ref{ch:results}. Four key observations
are examined in detail, and then the limitations and the ethical and legal
considerations are listed.

\section{The Gap Between Reducing Loss and Valid Generation}
\label{sec:disc-gap}

The most important finding of this project is a pattern that repeated in all seven
experimental generations of the direct track: the teacher-forced loss decreased
monotonically, while the autoregressive runtime condition was never met. Version two is
the clearest illustration of this pattern, because its best runtime behaviour occurs at
step 400, where the validation loss is \emph{higher}, and as the loss falls further until
step 2000, the runtime behaviour reaches zero (Figure~\ref{fig:v2}).

There are three possible explanations for this observation, and the project's evidence
helps to separate them.

\paragraph{Explanation one: a lack of learning capacity.} The low-rank adapter might lack
the capacity to represent Persian. Version v6.2 was designed specifically to falsify this
hypothesis: training and evaluation on the same 32 samples. The 32.20\% reduction of the
in-sample text loss shows that the adapter has the \emph{capacity} to fit the objective. This
explanation is therefore rejected.

\paragraph{Explanation two: exposure bias.} The explanation more consistent with the
evidence is the concept introduced in Section~\ref{sec:exposure-bias}. In training, the model
always receives the ``correct'' previous text and audio tokens as its condition, and
predicting one step ahead is an easy problem, especially in the ``inner monologue''
architecture, where the time-aligned text is present in the input. In autoregressive
generation, however, the model must build both the text and the audio from scratch. If the
first tokens drift slightly from the reference distribution, the model moves into a region it
has not seen in training, and errors grow cumulatively.

An observation that strengthens this explanation is visible in Table~\ref{tab:v45}: as the loss
decreases, the number of rows that produce \emph{any} speech falls from 9 to 4 or 5. That is, as
the model fits the objective better, it leans towards silence in independent generation. This
behaviour is not consistent with a model that has merely ``not learned enough''; it is
consistent with a model that has learned to predict correctly in the presence of the reference
input text and collapses in its absence.

Version v6.2 tried to compensate for exactly this: scheduled dropout of the text input raised
the share of hidden reference tokens from 0.25 to 0.75, so that the model would learn to rely on
its own predictions. This intervention kept the learning signal but did not transfer to valid
generation. A likely explanation is that with 32 training samples and 200 steps, the budget
available to relearn a stable generative behaviour is far below what is needed.

\paragraph{Explanation three: the English-centred base model.} The base model is strongly
English-centred. In the first generations it was observed that the output of intermediate
checkpoints is \emph{English} speech, not Persian. That is, the adapter could not shift the
distribution of the generated language and affected only acoustic details. Version four's
attempt to train the text embeddings selectively was a direct answer to this hypothesis, and it
was not enough either.

\paragraph{Methodological conclusion.} Regardless of which explanation contributes most, one
practical conclusion is certain, and in our view it is the most important scientific
contribution of this report:

\begin{quote}
In low-rank adaptation of a speech-to-speech model to a new language, the teacher-forced
validation loss is not a sufficient, or even a reliable, criterion for checkpoint selection. A
runtime test of autoregressive generation must be part of the selection criterion, not a check
after selection.
\end{quote}

Had the project based its selection criterion on loss alone, the step-8000 checkpoint of version
one, with the lowest validation loss and a significant improvement over the base model (1.7272
versus 3.5670 on held-out data), would have been reported as ``a Persian speech-to-speech
model'', while that same checkpoint produced almost silent audio and no text in the official
server.

\section{Why the Word Error Rate Is Much Higher Than the Character Error Rate}
\label{sec:disc-cer-wer}

In the re-recognition measure on the same 9 outputs of the mechanical panel with the
Qwen2.5-0.5B responder, the character error rate is 7.14\% and the word error rate 30.41\%; that
is, the word error rate is more than four times higher. At first sight this ratio looks
worrying, but it is to a large extent a \emph{definitional} effect in Persian.

Three mechanisms contribute to this gap. First, the \textbf{word boundary in Persian}. In the
definition used, the ZWNJ counts as a word boundary. If the recognition service writes a
compound word with a full space while the reference writes it with a ZWNJ, one word becomes two,
and this single difference produces two word edits, while at the character level there is only
one character of difference. Second, \textbf{inconsistent written forms}. Many Persian words
have several accepted written forms. Third, the \textbf{length effect}. The mean word length in
the data of this panel is about four characters; under these conditions, every one-character
error automatically becomes a complete word error.

There is supporting evidence for this interpretation in the data: in all 9 rows, the number of
characters in the hypothesis and the reference is almost the same (for example 52 and 52 in the
first row), while the number of words differs (11 and 12). That is, the character content is
almost preserved and the difference is mostly in segmentation.

Nevertheless, two limitations of this interpretation must be stated explicitly. First, part of
the error can come from wrong pronunciation by the speech engine rather than from
segmentation; separating these two shares is not possible without human listening. Second, the
recognition service takes part both in the system and in the measurement; their shared error can
shift the result in either direction. A character error rate of 7.14\% can therefore not be
translated into ``the generated speech is fully understandable for humans''.

\section{Why Adapting the Responder Lowered the Loss but Did Not Improve Quality}
\label{sec:disc-reference}

The experiment of Section~\ref{sec:res-lora} has an apparent contradiction: low-rank adaptation
of the responder reduced the development loss by 31.18\% but did not beat the baseline in the
semantic evaluation. The answer to this contradiction lies in the nature of the training
\emph{target}.

The training target of this adapter was to reproduce the ``next speaker's turn'' in the
conversational data. But the next speaker's turn in a natural conversation programme is not an
assistant's answer to a question. That turn can be a short affirmation, a change of topic, a
joke, a counter-question or the continuation of the speaker's own story. Trained on this target,
the model learns the \emph{style} of natural turn-taking better and the loss decreases, but it does
not learn instruction following and relevant answering, and may even damage them.

The direct and decisive evidence for this interpretation is in the final semantic test: the
``automatically aligned next-speaker reference'' arm received the lowest relevance score, 0.950,
lower even than the 0.5-billion-parameter baseline (1.425) (Figure~\ref{fig:semantic}). That is,
the automatic judge rated as irrelevant the very target the adapter had tried to imitate.

The practical consequence of this observation for later research is clear: automatically aligned
turns from natural conversation are \emph{valuable} material for learning speech timing,
turn-taking and interaction phenomena, but without a stage of selection and rewriting they are a
\emph{weak} target for training instruction-following answering. This point probably also
explains part of the failure of the direct track, because the direct model was trained on the
same pairs.

\section{A Structural Conflict Between the Requirements}
\label{sec:disc-conflict}

The analysis of Table~\ref{tab:positioning} reveals a structural conflict in the project
definition that is worth recording. Requirement R2 asks for full-duplex capability and
requirement R3 for a latency of at most 500 ms, while requirement R6 requires running on a GPU
with 12--24 GB of memory.

The models with \emph{native} full-duplex capability, which model the user's stream and
their own at the same time, did not fit within the target memory range at the time of this
project. The base model of this project's direct track used up to 28.640 GB of memory in
training. On the other hand, the first cascade architecture fitted within the target memory
range, but because the whole reply was built before playback, its latency (a median of 5584.7 ms
for full reply generation, Table~\ref{tab:cascade-rows}) was far from the target.

The project resolved this conflict by making the cascade streaming: speculative end-of-speech
detection, streaming token generation, clause chunking and synthesizing each clause as soon as
it is complete. The result was a mean of 372 and a 90th percentile of 431 ms in the browser
loopback test, with the GPU side at 11.1~GiB; full duplex was achieved with an external
two-stage controller instead of a native model. The cost of this solution is the loss of the
\emph{native} capabilities of full-duplex models, such as continuous turn prediction and the
assistant's own backchannels.

\section{The Value and Boundary of the Automatic Semantic Result}
\label{sec:disc-judge}

The result of the final semantic test is positive and passed all seven predefined conditions.
The design of this test has several methodological strengths: a qualification stage separate from
the final test, row indices and thresholds locked before running, blinding and randomization of
the arm order, deterministic decoding of the judge, and full validity of all 240 calls with a
repetition agreement rate of 1.00.

Nevertheless, three limitations separate it from a human result. First, the \textbf{same family
of judge and candidate}. The judge and the candidate are of the same model family; the literature
on automatic evaluation has shown that judges favour the output of models of their own
family~\cite{zheng2023mtbench}. Second, an \textbf{uncalibrated rubric}. The scoring rubric is not
calibrated against the scores of Persian speakers; we therefore do not know what level of real
user satisfaction a score of 3.150 out of 4 corresponds to. Third, the \textbf{semantic scope}. This
test measures only relevance and coherence and says nothing about the factual correctness of the
content, the safety of the reply or the naturalness of the speech.

Within these boundaries, the correct interpretation is: the cascade produces replies that an
automatic judge systematically, and by a wide margin (+1.725 in relevance), rates as more
relevant than the baseline. This is valid engineering evidence that the system works, not a
claim of perceptual quality.

\section{Limitations}
\label{sec:limitations}

The limitations of this project, stated explicitly and listed, are:

\begin{enumerate}
\item \textbf{Target hardware.} The evaluation was not run on a physical 12--24 GB GPU. The memory
      of the whole system was measured by limiting the language model server (11.1~GiB), but the
      compute speed is the H100's; latency on a consumer card will probably be higher.

\item \textbf{Automated loopback testing.} Latency and interruptions were measured with a headless
      browser and a fake microphone, not with a human user, a physical microphone and loudspeaker
      and real echo; in the simulation the echo was modelled as attenuated and delayed speech. In
      the final run all turns of the main panel are below 500 ms, but the maximum for simulated
      elderly speech (546) is above it, and with the browser's output latency 2\% of the turns of
      the main panel exceed 500.

\item \textbf{Human study.} Only a small listening study (5--10 participants, remote, based on
      recordings) was designed, and its results had not been received when this version was
      written. Live interaction with older adults was not measured.

\item \textbf{Automatic labels and judging.} The interruption labels come from extraction rules,
      Gemini and a list of listener tokens, reply quality was measured with an LLM judge and speech
      with an automatic recognizer; none of these is a human label. The two responder judges agree
      exactly only 44.6\% of the time.

\item \textbf{Interruption detection method.} The project definition asks for interruption detection
      with energy, pitch and MFCC features; a classifier based on these features alone reached only
      0.574 accuracy on real events, and the deployed controller relies on the speech recognizer,
      which brings the time to a complete stop to about half a second.

\item \textbf{Listening review of the data.} No row of the dataset is ``human-verified''.

\item \textbf{Direct track.} No experimental generation of Moshika and no Gemma speech-input adapter
      became eligible for deployment.

\item \textbf{Generalization.} The dataset is derived from four conversational channels, and the
      elderly user's voice in the tests is simulated (slowed machine speech), not recordings of older
      adults.

\item \textbf{Full reproduction.} The adapter tensors of versions two to five have been deleted, and
      only their hashes and results remain.
\end{enumerate}

\section{Ethical, Privacy and Legal Considerations}
\label{sec:ethics}

\subsection{Provenance and Rights of the Data}

The dataset is derived from an internal archive of Persian spoken content. Internal research use
of these resources was recorded with the supervisor's approval. This approval is not a ``licence to
redistribute'': the raw audio, the raw captions and their derivatives were not placed in the code
repository and are not published in the project's delivery package either. The statement ``publicly
available data has a free licence for training or redistribution'' is false and is not used in this
report.

\subsection{Privacy in the Human Study}

The system has three data retention modes: keeping the audio file, keeping only a reduced feature
vector, and keeping only metrics. The suggested default for the human study is the features-only
mode, so that no reconstructable speech of a participant is stored. An automatic test also checks
that no audio file is written in the metrics-only mode.

The principle of data minimization was also kept in the automatic evaluations: the result artifacts
do not store transcript text, reply text or judge explanations in plain text, and keep only hashes
and aggregate values. The retrospective descriptive analysis of Section~\ref{sec:res-cascade} was
also explicitly run without reading plain text.

\subsection{Licence and Its Boundaries}

The source code and documentation belonging to this project are published under the Apache-2.0
licence. This licence does \emph{not} cover the following, and the separate terms of each still
apply: third-party software components, the weights of language and speech models, the source
dataset and its derivatives, trained adapters, restricted media and captions, and the materials of
human-study participants. In particular, the speech engine used has a different licence, and its
obligations remain separate.

\subsection{Use and Misuse}

A system that generates natural Persian speech can be misused. In its current form, two factors
limit this risk: the output voice is a single, known synthesized voice and no voice cloning of
individuals is implemented; and the text responder works with a locked prompt. Nevertheless, the
project performed no content safety measurement, and the responder model can produce incorrect
content; using this prototype in any setting where an incorrect answer has serious consequences is
not permitted without a safety measurement layer.

\section{Statements This Report Avoids}
\label{sec:forbidden}

For transparency, the following eight statements \textbf{cannot be claimed} with the evidence of
this project. Each statement is quoted to make clear that these are not claims of this report but
an explicit list of things that \textbf{have not been said} and should not be concluded from this
project:

\begin{itemize}
\item ``Moshika was successfully fine-tuned for Persian.''
\item ``The system's latency is below 500 ms.''
\item ``The official interruption detection accuracy is above 80\%.''
\item ``The speech output is natural and fully understandable for humans.''
\item ``The system is suitable for use by older adults.''
\item ``The nine-row panel proves that the system generalizes to Persian speech.''
\item ``The source data have a free licence for redistribution or training.''
\item ``The dataset was labelled and verified by humans.''
\end{itemize}
